\documentclass[12pt,a4paper]{report}

\usepackage[spanish,es-nodecimaldot]{babel}
\usepackage[utf8]{inputenc}
\usepackage[T1]{fontenc}
\usepackage{lmodern}
\usepackage{amsmath,amssymb}
\usepackage{graphicx}
\usepackage{booktabs}
\usepackage{hyperref}
\usepackage{enumitem}
\usepackage{geometry}

\geometry{margin=2.5cm}

\title{Agente autónomo de ciberseguridad basado en aprendizaje por refuerzo}
\author{Javier Rivero Iglesias}
\date{Abril 2026}

\begin{document}

\maketitle
\tableofcontents

\chapter{Introducción}

Esta memoria de trabajo resume la evolución real del repositorio \texttt{TFG\_CYBER\_AI}.
El objetivo del proyecto es estudiar si un agente de aprendizaje por refuerzo puede
tomar decisiones binarias sobre tráfico de red:

\begin{itemize}
    \item \texttt{0 = PERMIT}
    \item \texttt{1 = BLOCK}
\end{itemize}

La historia del TFG tiene dos etapas claras. La primera fue una fase de validación
conceptual sobre NSL-KDD, útil para comprobar que el problema podía formularse como
una política binaria con recompensas sensibles al coste. La segunda, que es la línea
actualmente mantenida, migro a CICIDS2017, fijo un esquema canónico de 76 variables
de flujo y adopto QRDQN como algoritmo principal.

La memoria deja de presentar NSL-KDD como base final del sistema. Desde el punto de
vista del repositorio, la línea principal es ahora CICIDS2017 + QRDQN + esquema
canónico, mientras que NSL-KDD queda como antecedente histórico.

\chapter{Problema y alcance}

El trabajo se divide en dos fases:

\begin{enumerate}[label=\textbf{F\arabic*.}, leftmargin=*]
    \item \textbf{Fase 1: entrenamiento y validacion offline}. Se entrena un agente
    sobre datasets etiquetados y se mide su rendimiento con artefactos reproducibles.

    \item \textbf{Fase 2: inferencia offline en laboratorio}. Se aplican el modelo y
    el preprocesado aprendidos a flujos extraídos en un laboratorio privado para medir
    el efecto del cambio de distribución respecto al dataset.
\end{enumerate}

No forma parte de la línea base actual el bloqueo activo en producción. La Fase 2
sigue siendo de inferencia y análisis, no de despliegue defensivo autónomo.

\chapter{Evolución del repositorio}

\section{Etapa histórica: NSL-KDD + DQN}

La primera rama importante del proyecto utilizo NSL-KDD y agentes DQN junto a un
baseline supervisado Random Forest. Su función fue:

\begin{itemize}
    \item validar la formulación \texttt{PERMIT/BLOCK}
    \item explorar funciones de recompensa
    \item comparar RL frente a un baseline clásico
\end{itemize}

Los resultados históricos archivados muestran que Random Forest superaba
ligeramente a los primeros DQN:

\begin{center}
\begin{tabular}{llll}
\toprule
Experimento & Modelo & Accuracy & Recall ataque \\
\midrule
E01 & DQN & 0.7602 & 0.6000 \\
E02 & Random Forest & 0.7693 & 0.6150 \\
E05 & DQN & 0.7563 & 0.5955 \\
\bottomrule
\end{tabular}
\end{center}

Esta etapa fue útil como prueba de concepto, pero no quedo alineada con la futura
Fase 2 ni con una representación moderna de flujos.

\section{Migración a CICIDS2017 + esquema canónico + QRDQN}

La línea mantenida introdujo tres decisiones estructurales:

\begin{itemize}
    \item usar CICIDS2017 como dataset principal
    \item congelar un esquema canónico de 76 variables de flujo
    \item añadir una máscara de ausencia/presencia para construir observaciones de 152 dimensiones
\end{itemize}

La semántica de la máscara quedó fijada como:

\begin{itemize}
    \item \texttt{1 = present / valid}
    \item \texttt{0 = missing / imputed}
\end{itemize}

Sobre esa representación se consolidaron:

\begin{itemize}
    \item \texttt{src/canonical\_schema.py}
    \item \texttt{src/load\_cicids2017.py}
    \item \texttt{src/train\_rl\_defender.py}
    \item \texttt{src/validate\_checks.py}
    \item \texttt{src/validate\_leave\_one\_csv\_out.py}
    \item \texttt{scripts/predict\_real\_traffic\_v2.py}
\end{itemize}

\chapter{Línea base actual}

\section{Invariantes}

La implementación actual mantiene los siguientes invariantes:

\begin{itemize}
    \item \texttt{FEATURES\_CANON} contiene 76 características
    \item la observación final siempre tiene 152 dimensiones
    \item las etiquetas son \texttt{0 = BENIGN} y \texttt{1 = ATTACK}
    \item no se deben usar IPs, timestamps absolutos, Flow IDs ni proxies directos de puertos
\end{itemize}

\section{Recompensa}

El código actual está estandarizado en torno a:

\[
\text{reward\_config} =
\{ \texttt{tp}=1.5,\; \texttt{fp}=-1.5,\; \texttt{fn}=-5.0,\; \texttt{omission}=0.0 \}
\]

No obstante, varios artefactos históricos comprometidos fueron generados con otros
valores. El mejor run histórico de CICIDS2017, por ejemplo, uso \texttt{fp = -2.0}.
Por eso es necesario distinguir siempre entre \emph{defaults actuales} y
\emph{configuración histórica del run}.

\chapter{Resultados reproducibles}

\section{Entrenamiento CICIDS2017 + QRDQN}

Los principales runs comprometidos en el repositorio son:

\begin{center}
\begin{tabular}{llllllll}
\toprule
ID & Filas & Pasos & Split & Accuracy & Recall atk & F1 atk & Recompensa \\
\midrule
C01 smoke & 50k & 5k & random & 0.9697 & 0.99958 & 0.96922 & $1.5,-1.0,-5.0,0.0$ \\
C01 full & 250k & 100k & random & 0.99618 & 0.99980 & 0.99628 & $1.5,-1.0,-5.0,0.0$ \\
C02 fast & 100k & 10k & random & 0.9766 & 0.99959 & 0.98123 & $1.5,-1.0,-5.0,0.0$ \\
C03 full & 500k & 100k & random & 0.99859 & 0.99945 & 0.99876 & $1.5,-2.0,-5.0,0.0$ \\
\bottomrule
\end{tabular}
\end{center}

El mejor artefacto histórico comprometido es
\texttt{C03\_qrdqn\_cicids2017\_canonical\_full\_random\_20260223\_232439}, con:

\begin{itemize}
    \item accuracy: 0.99859
    \item precisión ataque: 0.99806
    \item recall ataque: 0.99945
    \item F1 ataque: 0.99876
\end{itemize}

\section{Validación}

La rama CICIDS2017 gano credibilidad cuando se pasó de medir solo \emph{random split}
a incorporar validaciones más exigentes:

\begin{center}
\begin{tabular}{llll}
\toprule
Check & Artefacto & Resultado clave & Lectura \\
\midrule
A & VAL\_checks\_A\_20260212\_235443 & accuracy 0.9939 & evaluacion directa fuerte \\
B & VAL\_checks\_B\_20260212\_235736 & shuffled 0.4773 & sin fuga aparente \\
C & VAL\_checks\_C\_20260213\_004847 & accuracy 0.84135 & generalizacion mas dura \\
\bottomrule
\end{tabular}
\end{center}

El Check C es especialmente importante porque entrena en los grupos
\texttt{Monday}, \texttt{Tuesday} y \texttt{Wednesday}, y evalúa sobre
\texttt{Thursday} y \texttt{Friday}. En ese régimen, el recall de ataque cae
a 0.52954, lo que deja claro que la generalización entre días sigue siendo
uno de los principales retos del proyecto.

La validación leave-one-CSV-out ya existe en código, pero no tiene todavía
un artefacto agregado comprometido bajo \texttt{runs/validation/}. Por tanto,
su metodología puede describirse, pero no sus métricas finales.

\chapter{Fase 2 y cambio de distribución}

La entrada mantenida para Fase 2 es \texttt{scripts/predict\_real\_traffic\_v2.py}.
Esta tubería reutiliza el modelo, el escalador y los percentiles del mejor run
de CICIDS2017 para inferir sobre flujos obtenidos en el laboratorio.

Los artefactos comprometidos muestran dos comportamientos muy distintos sobre
tráfico benigno:

\begin{center}
\begin{tabular}{lllll}
\toprule
Run & CSV & Block rate & Allow rate & z abs mean \\
\midrule
P2v2\_pred\_20260224\_004121 & flows\_benign.csv & 1.0 & 0.0 & 1.11489 \\
P2v2\_pred\_20260408\_230318 & flows\_benign.csv & 0.0 & 1.0 & 0.68709 \\
\bottomrule
\end{tabular}
\end{center}

La conclusión operativa es que la Fase 2 sigue siendo sensible al cambio de
distribución entre CICIDS2017 y tráfico real del laboratorio. Por eso cualquier
afirmación sobre comportamiento en Fase 2 debe ir siempre ligada al \texttt{RUN\_ID}
exacto y no formularse como una propiedad estable del sistema.

\chapter{Conclusiones y siguientes pasos}

El repositorio ya no es solo una colección de experimentos iniciales sobre NSL-KDD.
La línea mantenida basada en CICIDS2017, esquema canónico y QRDQN constituye una
base mucho mas sólida, con artefactos reproducibles y validaciones mejor definidas.

Sin embargo, permanecen tres riesgos principales:

\begin{itemize}
    \item la brecha entre \emph{random split} y generalización realista entre días
    \item la ausencia de un artefacto leave-one-CSV-out comprometido
    \item la sensibilidad de la Fase 2 al cambio de distribución
\end{itemize}

Los siguientes pasos razonables del TFG son:

\begin{itemize}
    \item conservar un run reproducible de leave-one-CSV-out
    \item repetir inferencia de Fase 2 sobre tráfico benigno y mixto con trazabilidad completa
    \item decidir si hace falta calibración adicional con datos del laboratorio
    \item mantener la memoria sincronizada con \texttt{docs/results.md} y los artefactos reales
\end{itemize}

\chapter{Fuentes internas del repositorio}

Esta memoria debe apoyarse en los siguientes documentos y artefactos internos:

\begin{itemize}
    \item \texttt{README.md}
    \item \texttt{.github/AGENT\_CONTEXT.md}
    \item \texttt{docs/results.md}
    \item \texttt{experiments/nslkdd\_experiments.md}
    \item \texttt{experiments/cicids2017\_qrdqn\_experiments.md}
    \item \texttt{runs/cicids2017/C03\_qrdqn\_cicids2017\_canonical\_full\_random\_20260223\_232439/}
    \item \texttt{runs/validation/VAL\_checks\_A\_20260212\_235443/}
    \item \texttt{runs/validation/VAL\_checks\_B\_20260212\_235736/}
    \item \texttt{runs/validation/VAL\_checks\_C\_20260213\_004847/}
    \item \texttt{runs/phase2/P2v2\_pred\_20260224\_004121/}
    \item \texttt{runs/phase2/P2v2\_pred\_20260408\_230318/}
\end{itemize}

\end{document}
